Computational Neuroscience is a rapidly developing research field that demands more and more computing resources for solving complex problems. Where in the past, models of signal processing in brain cells were strongly reduced -- due to the lack of computing power and high resolution experimental data -- they are now being developed in more detail. Emerging areas of research encompass large scale network simulations, three-dimensional simulations of neuronal signals (\cite{Xylouris2007, Xylouris2012, Breit2016, Breit2018, Breit2018b}) and multiscale modeling and simulation in neuroscience (\cite{Stepniewski2019b}).

In order to resolve biochemical and electrical signals of neurons, based on physical first principles, models defined by sets of highly nonlinear and coupled partial and ordinary differential equations need to be solved numerically and efficiently. Recently, biophysical models of electrical and biochemical signaling, in particular calcium dynamics, have been developed and implemented in the multiphysics simulation platform uG4 (\cite{Vogel2013}). uG4 has been shown to scale optimally on massively parallel architectures (\cite{Reiter12}). In various research projects the interplay between three-dimensional cell and organelle architecture and calcium signaling could be demonstrated (\cite{Grein14, Wittmann2009, Queisser11, Knodel2014, Breit2018, Breit2018b}). From these projects two currently active ones emerged that focus on two distinct cellular scales: (1) the ultrastructural dynamics at contact points between two neurons (synapses/synaptic spines) and (2) whole cell calcium dynamics including the calcium exchange mechanisms at the cellular plasma membrane and the endoplasmic reticulum. The importance of studying calcium dynamics are manifold. For instance, spine and cellular calcium dynamics are implicated in the long-term effects of repetitive transcranial magnetic stimulation (rTMS), a technique used in clinical settings to treat e.g. schizophrenia. At the whole cell level calcium-induced calcium release dynamics controlled by the endoplasmic reticulum have been shown critical in learning, memory formation, and cell survival (\cite{West2002, Milner1998, Bading2000, Hardingham1997}), as well as in neurodegenerative diseases, such as Alzheimer's and Parkinson's disease (\cite{Bezprozvanny2009, Chan2009, Cheung2008, Green2008}). 

The objectives in this research proposal are to enable high-throughput numerical simulations to study the effect of a multitude of geometric and biological parameters on calcium signaling pathways in neurons. An active NIH funded research project on multiscale modeling of rTMS, in collaboration with Dr. Opitz (Engineering, University of Minnesota), Dr. Jedlicka (Computational Neuroscience, University of Giessen, Germany), and Dr. Vlachos (Neurophysiology, University of Freiburg, Germany), as well as Ph.D. research focusses on the effect of synapse loss on synapse to cell nucleus communication will be supported by the computing resources applied for below. Resources will support the following two research objectives: 
\begin{enumerate}
\item At the level of synapses, synaptic spine morphology has been shown to change over time (\cite{Segal2014,Maggio2014,Verkhratsky2005,Finch2012,Jedlicka2017}). Along with this change, the intracellular cell organization adapts to ensure spine to dendrite coupling and calcium homeostasis in the spine head (\cite{Breit2018b}). In the proposed research, spine geometries that have been reconstructed in the rTMS project mentioned above, will be used to run 3D ultrastructural calcium simulations while scanning a parameter space consisting of receptor and calcium pump densities in the plasma- and endoplasmic membrane. These systematic simulation runs will identify critical parameter sets at which switching of neuronal behavior can be observed;{$\ $  in particular, we will determine under what physiological and geometric conditions we have stable calcium propagation from the spine-to-dendrite and wave propagation through the length of the dendrite}. Input for these simulations is voltage data derived from rTMS-induced electrical cell activity, which then drive {voltage dependent calcium channels which are} relevant for synaptic plasticity, {stable calcium wave propagation}, cell adaption, and survival. 

\item Whole cell calcium dynamics will be run by employing a 3D cell reconstruction method (\cite{Moerschel2017, Grein2020}) which uses 1D neuroanatomical data that is readily available from morphology databases such as NeuroMorpho.org (\cite{Ascoli2007}). This technique will allow high-throughput reconstruction of a large number of neurons on which parallel computations of cellular calcium signaling can be performed. Again, scanning a parameter space consisting of geometric parameters (e.g. ER size) and biophysical parameters, such as receptor densities, will ideally produce a rigorous understanding of critical states at which neurons loose their ability to reliably transmit calcium signals to the cell nucleus (which then triggers pathological states). 
\end{enumerate}
